\section{Completion and change of ground field}\label{sec:completion}

Smooth nonvanishing is now available in the order required by the
induction. We first finish the complex proof, including descent from
the good model to the original normal variety. We then descend the
globally generated Cartier multiple to an arbitrary algebraically
closed field of characteristic zero.

\begin{proof}[Proof of Theorem~\ref{thm:main} over \(\C\)]
Use dimension induction in the real-boundary category of
Proposition~\ref{prop:inductive-reduction}. In dimension zero the
variety is a point. Suppose the good-model assertion holds in all
smaller dimensions. The signed-representative argument
(Theorem~\ref{thm:signed}), the exclusions, and the jet estimates are
then available with exactly that lower-dimensional hypothesis.
Theorem~\ref{thm:smooth-nonvanishing} establishes smooth
nonvanishing in the present dimension. Proposition~\ref{prop:inductive-reduction}
therefore proves the real-boundary good-model assertion in the
present dimension. This is precisely the hypothesis needed to
continue to the next dimension.

Apply this to the rational lc pair \((X,B)\) of
Theorem~\ref{thm:main}. Let \(f:(X_d,B_d)\to(X,B)\) be a crepant dlt
model and let \((V,B_V)\) be its good log minimal model. Take a common
resolution \(a:W\to X_d\), \(q:W\to V\), and put \(p=f\circ a\).
Then
\[
 p^*(K_X+B)=a^*(K_{X_d}+B_d)=q^*(K_V+B_V).
\]
The first equality is crepancy. The second is the nef comparison in
Lemma~\ref{mmp:comparison}, since \(K_{X_d}+B_d=f^*(K_X+B)\) is nef.
The right side is semiample because \((V,B_V)\) is a good model, so
the pullback of \(K_X+B\) is semiample. Normality gives
\(p_*\OO_W=\OO_X\). Choose a sufficiently divisible Cartier multiple
\(m(K_X+B)\) whose pullback is globally generated. The projection
formula identifies its sections with those of its pullback.
Surjectivity of evaluation upstairs implies surjectivity downstairs:
otherwise a base point downstairs would make every pulled-back
section vanish on its nonempty fiber. Hence \(K_X+B\) is semiample.
\end{proof}

\begin{proposition}[Change of algebraically closed field]\label{prop:field-descent}
The conclusion of Theorem~\ref{thm:main} over \(\C\) implies its
conclusion over every algebraically closed field of characteristic zero.
\end{proposition}

\begin{proof}
Let \((X,B)\) be defined over such a field \(k\). Choose a finitely
generated subfield \(k_1\subset k\) over which the projective variety,
the rational boundary, a Cartier multiple of its adjoint, and a log
resolution are all defined. After enlarging \(k_1\) if necessary,
these data base change to the given data. Let \(k_0\) be its algebraic
closure inside \(k\), and denote the resulting pair by \((X_0,B_0)\).
Choose an embedding \(k_0\hookrightarrow\C\), and write
\((X_{\C},B_{\C})\) for the complex base change.

The chosen log resolution tests log canonicity after either
algebraically closed field extension: its boundary remains simple
normal crossing and all discrepancy coefficients remain unchanged.
The descended pair is consequently lc over both \(k_0\) and \(\C\).

Nefness is also invariant under such extensions. Indeed, a curve
over an extension is represented by a point of a relative Hilbert
scheme after its finite defining data have been spread over a
finite-type parameter scheme. The degree of a fixed line bundle is
constant in the resulting flat family. Specializing to a closed
point over the algebraically closed smaller field preserves a
negative degree, if one existed; some irreducible component of the
specialized curve would then have negative degree. This contradicts
nefness over that field. Conversely, curves over the smaller field
remain available after extension. Apply the converse to
\(k_0\subset k\): nefness of \(K_X+B\) gives nefness of
\(K_{X_0}+B_0\). Apply the forward direction to
\(k_0\hookrightarrow\C\): its complex base change is nef as well.

The complex case supplies an integer \(m>0\), enlarged to a multiple
of the Cartier index fixed over \(k_0\), for which the Cartier
line bundle \(M=\OO_{X_0}(m(K_{X_0}+B_0))\) becomes
globally generated over \(\C\). Proper flat base change for sections
identifies
\[
 H^0(X_0,M)\otimes_{k_0}\C
   = H^0(X_{\C},M_{\C}).
\]
Tensoring the evaluation map for \(M\) with \(\C\) therefore gives
the surjective complex evaluation map. Its coherent cokernel is
zero by faithful flatness. Extending from \(k_0\) to \(k\) preserves
this surjectivity, so the same integer \(m\) gives a globally
generated line bundle on the original \(X\).
\end{proof}

Together with Proposition~\ref{prop:field-descent}, the complex proof
establishes Theorem~\ref{thm:main} in its stated generality.
